\documentclass[border=12pt]{standalone}
\usepackage[T1]{fontenc}
\usepackage{lmodern}
\usepackage{amsmath,booktabs,array,multirow,makecell,xcolor}
\definecolor{Accent}{HTML}{315B7D}
\definecolor{Muted}{HTML}{5B6570}
\renewcommand{\arraystretch}{1.18}
\begin{document}
\begin{minipage}{16.2cm}
\raggedright
{\color{Accent}\bfseries Table [climate-associated cluster candidates]. Putative climate-associated LD clusters across BayPass configurations.}

\vspace{6pt}
\begin{center}\footnotesize
\begin{tabular}{@{}l l c r r r@{}}
\toprule
Population set & $\Omega$ (covariance) & Axis & \makecell{Candidates\\($\geq 5$)} & \makecell{Candidates\\($\geq 10$)} & \makecell{Median size\\(markers)} \\
\midrule
\multirow{4}{*}{\AA land-excluded (19 populations)} & \multirow{2}{*}{Internal (full data)} & PC1 & 14 & 8 & 88 \\
 &  & PC2 & 16 & 9 & 253 \\
\cmidrule(l){2-6}
 & \multirow{2}{*}{\makecell[l]{LD-pruned (fixed)\\ \textbf{[primary]}}} & PC1 & 9 & 3 & 126 \\
 &  & PC2 & 24 & 14 & 101 \\
\midrule
\multirow{4}{*}{Full (20 populations)} & \multirow{2}{*}{Internal (full data)} & PC1 & 7 & 6 & 497 \\
 &  & PC2 & 44 & 28 & 98 \\
\cmidrule(l){2-6}
 & \multirow{2}{*}{LD-pruned (fixed)} & PC1 & 4 & 3 & 449 \\
 &  & PC2 & 16 & 11 & 260 \\
\bottomrule
\end{tabular}
\end{center}

\vspace{2pt}{\scriptsize\color{Muted} A cluster is a putative candidate when at least 5 (or 10) of its member markers reach $\mathrm{BF(dB)}\geq 15$. Genome-wide median cluster size $=9$ markers (32,854 clusters with a consensus genotype); the excess size of candidate clusters over this median reflects the size dependence of the member-count criterion. The eMLG gate restricts this set to clusters that track climate as a unit (per-cluster table).\par}
\end{minipage}
\end{document}
